\hypertarget{appendix-eb336149c00a}{}
\section*{Floating-point infinities}
\section{Floating-point infinities}\label{floating-point-infinities}

\textbf{Curriculum depth:} Introduced · \textbf{First seen:} C004 · \textbf{Latest development:} C004

\subsection{Definition and mechanism}\label{definition-and-mechanism}

All-one exponent and zero fraction encode signed infinity. This is not maximum finite. Infinity minus itself is NaN; overflow's result also depends on rounding direction.

\subsection{Prerequisites}\label{prerequisites}

\begin{itemize}
\tightlist
\item
  \hyperref[appendix-5057f9ad327f]{Binary32} --- The concrete fields provide the representation used in this mechanism.
\end{itemize}

\subsection{Used by}\label{used-by}

\begin{itemize}
\tightlist
\item
  \hyperref[appendix-0ec599f8ba27]{Numerical tolerance policies} --- uses this mechanism in its explanation.
\item
  \hyperref[appendix-7f1e17dcafb1]{Robust norm and intermediate range} --- uses this mechanism in its explanation.
\end{itemize}

\subsection{Recall and next depth}\label{recall-and-next-depth}

Reconstruct the concrete example in the definition and identify its validity assumptions.

\textbf{Pending:} Define non-finite input policy before applying a relative tolerance formula.

\subsection{Source and navigation}\label{source-and-navigation}

\href{Continuous\%20Notes/Computing\%20Foundations\%20-\%20From\%20Bits\%20to\%20C++\%20and\%20Systems\%20APIs\#A\%20stretching\%20ruler\%20for\%20real-valued\%20quantities}{M004 teaching chapter} · \hyperref[appendix-fd090d982a15]{Per-part definitions} · \href{}{Dependency Map (02-Foundations-II, appendix-4ad47bdd269a)}

\subsection{Reconstruction depth}\label{reconstruction-depth}

IEEE-754 infinities represent values beyond the finite exponent range and results such as nonzero finite division by signed zero. They participate in arithmetic according to defined extended-real-like rules. Adding a finite value to positive infinity stays positive infinity, while indeterminate forms such as infinity minus infinity produce NaN.

Infinity can be useful as a sentinel in algorithms, such as an initial shortest-path distance, because comparisons with finite positive values behave naturally. But it is not automatically a valid physical or engine state. A transform position, mass, texture coordinate or timestamp becoming infinite usually signals overflow or invalid input and should be caught at a subsystem boundary.

A robust numeric API distinguishes ``the floating type can encode infinity'' from ``this domain accepts infinity.'' Tests should exercise overflow, division by zero under the active floating environment, comparisons, finite-value validation and serialization behavior. Fast-math modes may weaken exceptional-value assumptions, so production guarantees must match compile configuration.
